% Course extension (not an official 2023 A scored page).
% Numbers from ../results/bioeconomic_pareto.csv,
% bioeconomic_tradeoff.csv, monthly_metrics.csv, impact_factors.csv.
% Compile: pdflatex hoa_community_statement.tex
\documentclass[11pt]{article}
\usepackage[margin=1in]{geometry}
\usepackage[hidelinks]{hyperref}
\usepackage{setspace}
\pagestyle{empty}
\setstretch{1.08}

\begin{document}

\begin{center}
{\Large\bfseries Rethinking the Yellow Bloom:}\\[0.25em]
{\large A Science-Based Coexistence Guide for Sustainable Turf Management}\\[0.6em]
{\small Homeowners Association / Parks Department briefing
\quad GoHiMCM 2023~A study team}
\end{center}

\vspace{0.6em}
\noindent
Dear Board and Parks staff,

The yellow week in April is not a failure of the crew.
It is the first bloom of a plant that, on our temperate weekly simulation
of \textbf{one puffball just west of a one-hectare common}, occupies
$\mathbf{45.11\%}$ of cells by December if nobody mows
(\texttt{monthly\_metrics.csv}).
That is a west-in patch with $4674$ flowering adults and a much larger
seedling count---not a uniform carpet, and not the daily classroom CA
that starts a plant in the dead center and fills $100\%$ of the grid.

Against kudzu and Japanese knotweed, scored on the same six-leaf sheet,
dandelion's impact is $0.151$ versus $0.933$ and $0.570$
(\texttt{impact\_factors.csv}).
That ranking is only inside those three rows.
It is a reason to stop panic-spraying the whole common for one fence-line
patch. It is not a reason to ignore seed rain.

\section*{The two-zone, two-phase rule}

Do not run one calendar on every square meter.

\paragraph{Zone A --- early-spring bee verge.}
Mark a strip along the west fence (and any already-yellow median)
as a \emph{pollinator verge} from March through the first two weeks of
open bloom.
Leave those heads standing.
In our daily window model the first real bloom is April;
keeping about fourteen days of flowers is the entire pollinator term
$B_{\mathrm{pollinator}}$.

\paragraph{Zone B --- show turf.}
Play fields, front walks, and HOA photo lawns stay on a height rule,
not a herbicide calendar.

\paragraph{Phase 1 --- keep the bloom.}
For Zone~A, the two weeks before and into first bloom are a no-cut
window. That is the ``retain for bees'' half of Strategy~4.

\paragraph{Phase 2 --- cut 72 hours before seed fly.}
When the florets close and the pappus clock starts, cut Zone~A once,
low deck, about \textbf{three days (72 hours) before the seed heads open}.
In the model that is a single spring cut after $\approx 11$ days of open
flower. It is cheap: one mechanical pass, not ten 2,4-D rounds.

\section*{What the numbers actually say}

The spring window is not a silver bullet.
On the daily CA, Strategy~4 keeps $\mathbf{51.8\%}$ of the wild
early-spring pollinator score (not $84.5\%$) and only
$\mathbf{4.4\%}$ fewer year-end plants than doing nothing
(\texttt{bioeconomic\_tradeoff.csv}).
Autumn bloom refills the hectare: year-end cover is still $100\%$
in that geometry.
So Phase~2 must be followed by a \emph{turf rule} on Zone~B.

On the weekly cost--cover front
(\texttt{bioeconomic\_pareto.csv}),
mowing every two weeks at full height reduction reaches cover
$\mathbf{0.0064}$ at relative cost $39$.
Mowing every week costs $78$ for the same cover---wasted labor.
Never-mow is cover $0.4511$ at cost $0$.
If the Board wants bees \emph{and} a tidy photo lawn:
keep Zone~A through Phase~1--2, then put Zone~B on the biweekly cadence
for the rest of bloom season (including September--October).
Do not spray the hectare because one verge went yellow.

This is the \emph{temperature--bloom dual-response} card we distilled from
the PPO run (\texttt{rl\_policy\_evaluation.csv}): the network's greedy
year is ``do nothing'' (net $-130$, same as rewilding; 14-day scalping
nets $-203$).
So the letter is not ``the neural net mows at 6.8\% cover.''
It is: \textbf{when air is still in the March--April bloom, do not cut
Zone~A}; \textbf{when you need a photo lawn after seed-set, use the
biweekly rule}, not weekly poison.

Drought already cuts weekly month-12 cover from $0.4511$ to $0.2046$.
After a dry summer, revisit the cadence before adding chemistry.

\section*{What we are asking you to print}

\begin{enumerate}
  \item Stake Zone~A in March. No cut until Phase~2.
  \item One low-deck cut on Zone~A 72 hours before seed fly.
  \item Zone~B: biweekly mow through spring and fall bloom; skip weekly
        scalping.
  \item Herbicide only on documented knotweed or kudzu, not on
        \emph{Taraxacum}.
\end{enumerate}

\vspace{1.2em}
\noindent
Respectfully,\\[0.6em]
GoHiMCM 2023~A study team\\
{\small Sources: \texttt{results/monthly\_metrics.csv},
\texttt{bioeconomic\_pareto.csv},
\texttt{bioeconomic\_tradeoff.csv},
\texttt{impact\_factors.csv},
\texttt{rl\_policy\_evaluation.csv}.}

\end{document}
